\section{Discussion and Limitations}\label{sec:discussion}
\subsection{When does RL help link adaptation?}
In this study RL helped only where the decision had structure that a static offset cannot use. On full-buffer goodput (M3, M4) every useful input is already summarised by a calibrated offset per user class, and the residual policy learned to leave its base rule alone at RTT~\NCfgRtt{} and did worse than it at RTT~\lit{0}. With deadline traffic and delayed HARQ feedback (M1) the policy found something to add. Its corrections are most frequent when HARQ processes are busy, when little of the deadline budget remains and during interference bursts (Section~\ref{sec:behaviour}). This is consistent with the policy trading block-error risk against the chance of finishing a packet in time, a trade-off that a table indexed by a few bins approximates only coarsely. We did not test this mechanism.

Three qualifications bound the positive result. First, the margin over the strongest calibrated rule is small: \NGsixCtxJMar\,\%, or \NGsixAbsPpCtxJ{} percentage points of on-time delivery. Second, it is not robust across training seeds: only \NGsixCtxJAbove{} of \NCfgSeeds{} seeds beat the combined table, and the hierarchical interval [\NGsixCtxJHierLo, \NGsixCtxJHierHi]\,\% includes zero. The ten policies also act very differently (the share of upward corrections ranges from \NGoneAllSeedMin{} to \NGoneAllSeedMax\,\% across seeds). A practitioner who trains a single policy has roughly even odds of beating the combined table: the probability of improvement is \NGsixCtxJPoI{} ([\NGsixCtxJPoILo, \NGsixCtxJPoIHi]). Third, the gain is not uniform across users. Against the v17 Context table, users above \NAoneEdgeLo~m/s gain (\NAoneMidMar\,\% and \NAoneFastMar\,\%), while users below it have a margin of \NAoneSlowMar\,\% (hierarchical \NLevHier\,\% CI [\NAoneSlowLo, \NAoneSlowHi]\,\%), and these slow users make up \NAoneSlowShare\,\% of user-traces. The policy's actions barely depend on speed, so a speed-aware residual or a speed-gated deployment might remove the loss; we have not tried either.

Table~\ref{tab:summary} collects where the learned policy helped in this study, and by how much.

\begin{table*}[t]
\caption{Summary: where the residual A2C policy helped. Margins are trace-level (Rule~R1); the seed-level column gives the hierarchical seed$\times$trace \NLevHier\,\% CI against the same comparator where it was computed (Rule~R2: robust if the lower bound is above zero). $^\dagger$Against the re-tuned OLLA of G6.}\label{tab:summary}
\centering\footnotesize
\setlength{\tabcolsep}{3pt}
\begin{tabular}{@{}lp{5.2cm}cp{3.0cm}p{3.6cm}@{}}
\toprule
Setting & Comparator & Margin (\%) & Beats (R1) & Seed level: hier.\ \NLevHier\,\% CI (\%)\\
\midrule
M1: deadline, RTT~\NCfgRtt & OLLA, cold start / after warm-up & \NMonePrimOLLAMar{} / \NWoneOLLAWarmMar & \NMonePrimOLLABeats{} / \NWoneOLLAWarmBeats & [\NGsixOLLATDHierLo, \NGsixOLLATDHierHi]$^\dagger$\\
M1 & best v17 rule (Context table, \NCfgCalibCtx{} traces) & \NMonePrimCtxMar & \NMonePrimCtxBeats & [\NGMoneCtxHierLo, \NGMoneCtxHierHi]\\
M1 & best single-cue rule, \NGsixCalTraces{} traces (Context-500) & \NGsixCtxDMar & \NGsixCtxDBeats & [\NGsixCtxDHierLo, \NGsixCtxDHierHi]\\
M1 & \textbf{combined table (Context-joint-500)} & \NGsixCtxJMar & \NGsixCtxJBeats{} (not after project-wide correction) & \textbf{[\NGsixCtxJHierLo, \NGsixCtxJHierHi]}\\
M1, users $<\NAoneEdgeLo$~m/s & v17 Context table & \NAoneSlowMar & -- & [\NAoneSlowLo, \NAoneSlowHi]\\
M1, non-oracle base (NR1) & best v17 rule & \NNRoneCtxMar & yes & not stored (NR2 pending)\\
M2: deadline, RTT~\lit{0} & best v17 rule & \NMtwoPrimCtxMar & \NMtwoPrimCtxBeats & [\NGMtwoCtxHierLo, \NGMtwoCtxHierHi]\\
M3: goodput, RTT~\NCfgRtt & Fixed & \NMthreePrimFixedMar & \NMthreePrimFixedBeats & --\\
M4: goodput, RTT~\lit{0} & Fixed & \NMfourPrimFixedMar & \NMfourPrimFixedBeats & --\\
\bottomrule
\end{tabular}
\end{table*}

\subsection{Calibrate first}
The largest effect in this study is not the learner but the comparator. Against OLLA started cold the learner wins by \NWoneOLLAColdMar\,\%; a warm-up of \NWoneTTI{} TTIs reduces this to \NWoneOLLAWarmMar\,\%. One calibrated offset (Fixed) captures \NMoneCalibShareFixed\,\% of the learner's gain over OLLA, and the v17 Context table captures \NMoneCalibShareCtx\,\%. Exhaustive re-calibration on \NGsixCalTraces{} traces and a combined table remove a further \NGsixClosure\,\% of the remaining margin. Each step changes the conclusion a reader would draw. The steps are not all cheap. Calibrating the single-cue rules on \NGsixCalTraces{} traces took \NGsixSingleHours~h of simulation, but the coordinate-ascent search for the combined table took \NGsixSearchHours~h, against \NMoneHours~h to train and evaluate the ten M1 policies. A fair comparison therefore needs a comparable tuning budget for the rules, and that budget should be reported, as for the learner~\cite{ferraridacrema2019progress}. We therefore read the M1 result as: \emph{calibrate first}; at the studied point the learned residual then adds a fraction of a percentage point on average, with a seed-dependent sign against the best table.

The calibration study also says something about the base rule. On the evaluation traces the non-oracle Fixed-500 offset scores \NGsixBaseGapPct\,\% above the class-oracle FixedPU-500 table, although FixedPU was the base rule selected for M1 on \NCfgCalib{} calibration traces. A base rule picked on a small calibration set can thus be the weaker one even when it has more information. This is consistent with NR1, where the residual policy on the Fixed base was \NNRoneMoneMar\,\% above the M1 policy on the FixedPU base. NR1 used M1's evaluation traces and has no seed-level analysis, so it does not establish a better configuration. The pre-registered replication NR2 (fresh traces, fresh seeds, the frozen combined table as primary comparator) decides that question.

\subsection{Cost}
The actor has \NParamsActor{} parameters and \NMacsActor{} multiply-accumulates per decision. On an \NAthreeCpu{} laptop CPU a decision takes \NAthreeLoop~$\mu$s in a batched loop, and a full TTI for \NCfgUsers{} users \NAthreeTti~$\mu$s, \NAthreeRatio{} times the cost of the FixedPU rule in the same simulator. This is small. Training is not negligible (\NCfgEpDl{} episodes per seed, ten seeds), but in this study the search for the strongest comparator cost more simulation time than training the learner (Section~\ref{sec:discussion}-B).

\subsection{Limitations}
\begin{enumerate}
\item \emph{Simulator.} One macro cell, $N=\NCfgUsers$ users and one user scheduled per TTI. There is no multi-cell interference coordination, carrier aggregation or frequency-selective scheduling. Block errors use an SINR-to-BLER abstraction rather than NR LDPC link-level tables, and there is no cross-validation in ns-3 or Sionna. All results come from one MATLAB code base.
\item \emph{Seed robustness.} The margin over the combined table is not robust to seed resampling (Rule~R2). NR1 cannot be assessed at seed level because per-seed values were not stored.
\item \emph{Heterogeneous gain.} For users below \NAoneEdgeLo~m/s the margin against the v17 Context table is \NAoneSlowMar\,\%. The A1 breakdown was not repeated against the G6 tables.
\item \emph{Operating point.} \wordingsel{G10}{G10}{verdict}{unrun}
\item \emph{Oracle information.} The M1 base rule FixedPU uses the true speed class, which the gNB does not observe, and the combined table uses it too. The learner's own state is gNB-observable. NR1 removes the oracle from the base rule but reuses M1's evaluation traces; NR2 is pending.
\item \emph{Scoring window.} Traces score \NCfgScored{} TTIs after a burn-in. OLLA starts cold unless warmed up (W1).
\item \emph{Selection over versions.} The residual recipe was chosen over \NVersions{} versions of the campaign. M1 is a confirmatory test of that recipe on fresh data, and multiplicity is accounted for over the registry, but selection effects cannot be fully excluded. The registry counts in this draft predate the v17x robustness stages, whose rows are still to be added.
\item \emph{Definitions.} Conclusions about multi-process HARQ depend on the PF-averaging definition (Section~\ref{sec:pitfalls}). We use the scheduling-time definition throughout.
\item \emph{Learned baselines.} No learned LA scheme from the literature (for example LOLLA or SALAD) was re-implemented in our simulator.
\end{enumerate}
